% !TEX program = xelatex
\documentclass[12pt,a4paper]{ctexbook}

% ==================== 页面 ====================
\usepackage[top=2.5cm,bottom=2.5cm,left=2.5cm,right=2.5cm]{geometry}

% ==================== 字体 ====================
\usepackage{fontspec}
\usepackage{iffont}
\IfFontExistsTF{Consolas}
  {\setmonofont{Consolas}}
  {\setmonofont{DejaVu Sans Mono}[Scale=MatchLowercase]}

% ==================== 颜色 ====================
\usepackage[dvipsnames,svgnames,x11names]{xcolor}
\definecolor{codebg}{HTML}{F5F5F5}
\definecolor{codeframe}{HTML}{E0E0E0}
\definecolor{cppkeyword}{HTML}{1A1AFF}
\definecolor{cppcomment}{HTML}{2E8B2E}
\definecolor{cppstring}{HTML}{B22222}
\definecolor{linenumgray}{HTML}{999999}
\definecolor{csharpkeyword}{HTML}{8B008B}
\definecolor{csharptype}{HTML}{006400}
\definecolor{cmakekw}{HTML}{0050A0}
\definecolor{cmakeheadbg}{HTML}{2E5090}
\definecolor{xmakeheadbg}{HTML}{5B3A8C}
\definecolor{vsheadbg}{HTML}{0078D4}
\definecolor{xcodeheadbg}{HTML}{1470B0}
\definecolor{makeheadbg}{HTML}{6B8E23}
\definecolor{androidheadbg}{HTML}{3DDC84}
\definecolor{oeheadbg}{HTML}{C71A36}
\definecolor{cmakekwcolor}{HTML}{003399}
\definecolor{xmakekwcolor}{HTML}{6A1B9A}
\definecolor{bashkw}{HTML}{1A1AFF}
\definecolor{cfgheadbg}{HTML}{2E4057}
\definecolor{algheadbg}{HTML}{1A3C5E}

% ==================== listings ====================
\usepackage{listings}

\lstdefinelanguage{CMake}{
  morekeywords={cmake\_minimum\_required,project,set,option,if,elseif,else,endif,
    foreach,endforeach,while,endwhile,function,endfunction,macro,endmacro,
    add\_executable,add\_library,add\_subdirectory,add\_custom\_target,add\_custom\_command,
    target\_link\_libraries,target\_include\_directories,target\_compile\_features,
    target\_compile\_definitions,target\_compile\_options,target\_sources,
    find\_package,find\_path,find\_library,find\_program,
    include,include\_directories,link\_directories,
    install,export,configure\_file,file,
    set\_target\_properties,set\_property,get\_property,
    message,execute\_process,enable\_testing,add\_test,
    include\_external\_msproject,add\_dependencies,
    cmake\_parse\_arguments,list,string,math,
    set\_source\_files\_properties,set\_directory\_properties,
    generator,cmake\_policy,cmake\_path,fetchcontent\_declare,
    fetchcontent\_makeavailable,fetchcontent\_getproperties,
    cpack\_add\_component,include\_cpack,ctest},
  sensitive=false,
  morecomment=[l]{\#},
  morestring=[b]",
}

\lstdefinelanguage{XMake}{
  morekeywords={add\_requires,target,add\_files,set\_kind,set\_languages,
    set\_warnings,set\_optimize,add\_includedirs,add\_linkdirs,add\_links,
    add\_defines,add\_cxxflags,add\_cflags,add\_shflags,add\_ldflags,
    set\_targetdir,set\_basename,set\_extension,
    add\_rules,add\_deps,add\_packages,add\_imports,
    on\_build,on\_clean,on\_install,on\_run,on\_package,
    before\_build,after\_build,before\_clean,after\_clean,
    set\_toolchains,set\_plat,set\_arch,
    add\_requireconfs,add\_repositories,
    option,add\_options,set\_showmenu,set\_description,
    task,set\_menu,on\_run,includes},
  sensitive=true,
  morecomment=[l]{--},
  morestring=[b]",
  morestring=[b]',
}

\lstdefinelanguage{Bash}{
  morekeywords={cmake,xmake,make,ninja,clang,clang++,gcc,g++,cl,
    export,cd,mkdir,rm,cp,mv,echo,cat,grep,sed,awk,find,xargs,
    curl,wget,sudo,apt,yum,brew,choco,winget,
    powershell,pwsh,dotnet,nuget,bitbake,repo,source,
    ./configure,./autogen.sh,autoreconf,meson,
    ndk-build,gradle,./gradlew,source,inherit},
  sensitive=true,
  morecomment=[l]{\#},
  morestring=[b]",
  morestring=[b]',
}

\lstdefinelanguage{MSBuild}{
  morekeywords={Project,PropertyGroup,ItemGroup,Target,Import,
    ItemDefinitionGroup,ClCompile,ClInclude,Link,ResourceCompile,
    Configuration,Platform,OutputDirectory,IntermediateDirectory,
    PreprocessorDefinitions,AdditionalIncludeDirectories,
    AdditionalLibraryDirectories,AdditionalDependencies,
    SubSystem,TargetMachine,GenerateDebugInformation,
    ImportGroup,PropertySheet,ConfigurationType,
    PlatformToolset,CharacterSet,WholeProgramOptimization,
    RuntimeLibrary,WarningLevel,Optimization,
    ItemDefinitionGroup,PreBuildEvent,PostBuildEvent,CustomBuildStep},
  sensitive=true,
  morecomment=[l]{<!--},
  morecomment=[l]{-->},
  morestring=[b]",
}

\lstdefinelanguage{BitBake}{
  morekeywords={inherit,require,include,addtask,deltask,addhandler,
    export,unset,INHERIT,PACKAGECONFIG,EXTRA\_OECONF,
    do\_configure,do\_compile,do\_install,do\_fetch,do\_unpack,
    do\_patch,do\_package,do\_populate\_sysroot,
    S,B,D,WORKDIR,TMPDIR,PV,PR,PN,PE,
    SRC\_URI,SRCREV,S,DEPENDS,RDEPENDS,RRECOMMENDS,
    LICENSE,LIC\_FILES\_CHKSUM,
    BBCLASSEXTEND,PROVIDES,ALTERNATIVE,
    IMAGE\_INSTALL,IMAGE\_FEATURES,
    PACKAGE\_ARCH,TUNE\_FEATURES,MACHINE},
  sensitive=true,
  morecomment=[l]{\#},
  morestring=[b]",
  morestring=[b]',
}

\lstdefinelanguage{Plist}{
  morekeywords={plist,dict,key,string,array,integer,true,false,
    CFBundleDevelopmentRegion,CFBundleExecutable,CFBundleIdentifier,
    CFBundleName,CFBundlePackageType,CFBundleShortVersionString,
    CFBundleVersion,NSMainStoryboardFile},
  sensitive=true,
  morecomment=[l]{<!--},
  morecomment=[l]{-->},
  morestring=[b]",
}

% ---- listing styles ----
\lstdefinestyle{cppstyle}{
  language=C++,
  basicstyle=\small\ttfamily,
  keywordstyle=\color{cppkeyword}\bfseries,
  commentstyle=\color{cppcomment}\itshape,
  stringstyle=\color{cppstring},
  numberstyle=\tiny\color{linenumgray},
  numbers=left,
  frame=single,
  rulecolor=\color{codeframe},
  framesep=6pt,
  breaklines=true,
  tabsize=4,
  columns=flexible,
  keepspaces=true,
  backgroundcolor=\color{codebg},
  showstringspaces=false,
  morekeywords={constexpr,consteval,constinit,decltype,auto,
    concept,requires,co_await,co_yield,co_return,
    noexcept,override,final,nullptr,static_assert,
    template,typename,namespace,using,mutable,volatile,
    explicit,operator,friend,inline,virtual,
    thread_local,alignas,alignof,
    if,consteval,constexpr},
  deletekeywords={int,char,bool,float,double,long,short,unsigned,signed,void,wchar_t},
  emph={size_t,int32_t,uint32_t,int64_t,uint64_t,string,vector,map,set,
    unique_ptr,shared_ptr,optional,variant,any,span,string_view},
  emphstyle=\color{teal!80!black},
  escapeinside={(*@}{@*)},
}
\lstdefinestyle{cstyle}{
  language=C,
  basicstyle=\small\ttfamily,
  keywordstyle=\color{cmakekw}\bfseries,
  commentstyle=\color{cppcomment}\itshape,
  stringstyle=\color{cppstring},
  numberstyle=\tiny\color{linenumgray},
  numbers=left,
  frame=single,
  rulecolor=\color{codeframe},
  framesep=6pt,
  breaklines=true,
  tabsize=4,
  columns=flexible,
  keepspaces=true,
  backgroundcolor=\color{codebg},
  showstringspaces=false,
  deletekeywords={int,char,bool,float,double,long,short,unsigned,signed,void,wchar_t},
  escapeinside={(*@}{@*)},
}
\lstdefinestyle{cmakestyle}{
  language=CMake,
  basicstyle=\small\ttfamily,
  keywordstyle=\color{cmakekwcolor}\bfseries,
  commentstyle=\color{cppcomment}\itshape,
  stringstyle=\color{cppstring},
  numberstyle=\tiny\color{linenumgray},
  numbers=left,
  frame=single,
  rulecolor=\color{codeframe},
  framesep=6pt,
  breaklines=true,
  tabsize=4,
  columns=flexible,
  keepspaces=true,
  backgroundcolor=\color{codebg},
  showstringspaces=false,
  escapeinside={(*@}{@*)},
}
\lstdefinestyle{xmakestyle}{
  language=XMake,
  basicstyle=\small\ttfamily,
  keywordstyle=\color{xmakekwcolor}\bfseries,
  commentstyle=\color{cppcomment}\itshape,
  stringstyle=\color{cppstring},
  numberstyle=\tiny\color{linenumgray},
  numbers=left,
  frame=single,
  rulecolor=\color{codeframe},
  framesep=6pt,
  breaklines=true,
  tabsize=4,
  columns=flexible,
  keepspaces=true,
  backgroundcolor=\color{codebg},
  showstringspaces=false,
  escapeinside={(*@}{@*)},
}
\lstdefinestyle{bashstyle}{
  language=Bash,
  basicstyle=\small\ttfamily,
  keywordstyle=\color{bashkw}\bfseries,
  commentstyle=\color{cppcomment}\itshape,
  stringstyle=\color{cppstring},
  numberstyle=\tiny\color{linenumgray},
  numbers=left,
  frame=single,
  rulecolor=\color{codeframe},
  framesep=6pt,
  breaklines=true,
  tabsize=4,
  columns=flexible,
  keepspaces=true,
  backgroundcolor=\color{codebg},
  showstringspaces=false,
  escapeinside={(*@}{@*)},
}
\lstdefinestyle{xmlstyle}{
  language=MSBuild,
  basicstyle=\small\ttfamily,
  keywordstyle=\color{vsheadbg}\bfseries,
  commentstyle=\color{cppcomment}\itshape,
  stringstyle=\color{cppstring},
  numberstyle=\tiny\color{linenumgray},
  numbers=left,
  frame=single,
  rulecolor=\color{codeframe},
  framesep=6pt,
  breaklines=true,
  tabsize=4,
  columns=flexible,
  keepspaces=true,
  backgroundcolor=\color{codebg},
  showstringspaces=false,
  escapeinside={(*@}{@*)},
}
\lstdefinestyle{bbstyle}{
  language=BitBake,
  basicstyle=\small\ttfamily,
  keywordstyle=\color{oeheadbg}\bfseries,
  commentstyle=\color{cppcomment}\itshape,
  stringstyle=\color{cppstring},
  numberstyle=\tiny\color{linenumgray},
  numbers=left,
  frame=single,
  rulecolor=\color{codeframe},
  framesep=6pt,
  breaklines=true,
  tabsize=4,
  columns=flexible,
  keepspaces=true,
  backgroundcolor=\color{codebg},
  showstringspaces=false,
  escapeinside={(*@}{@*)},
}
\lstset{style=cmakestyle}

% ==================== tcolorbox ====================
\usepackage[most]{tcolorbox}

\newtcolorbox{guideline}[1][]{
  enhanced,breakable,
  colback=blue!5!white,colframe=blue!75!black,
  fonttitle=\bfseries,
  title=配置要点,#1,
  boxed title style={colback=blue!75!black},
  attach boxed title to top left={yshift=-2mm,xshift=2mm},
  arc=2mm,boxrule=0.8mm,
}
\newtcolorbox{compare}[1][]{
  enhanced,breakable,
  colback=green!3!white,colframe=green!50!black,
  fonttitle=\bfseries,
  title=构建系统对比,#1,
  boxed title style={colback=green!50!black},
  attach boxed title to top left={yshift=-2mm,xshift=2mm},
  arc=2mm,boxrule=0.8mm,
}
\newtcolorbox{warning}[1][]{
  enhanced,breakable,
  colback=red!5!white,colframe=red!75!black,
  fonttitle=\bfseries,
  title=常见陷阱,#1,
  boxed title style={colback=red!75!black},
  attach boxed title to top left={yshift=-2mm,xshift=2mm},
  arc=2mm,boxrule=0.8mm,
}
\newtcolorbox{note}[1][]{
  enhanced,breakable,
  colback=orange!5!white,colframe=orange!80!black,
  fonttitle=\bfseries,
  title=注意,#1,
  boxed title style={colback=orange!80!black},
  attach boxed title to top left={yshift=-2mm,xshift=2mm},
  arc=2mm,boxrule=0.8mm,
}
\newtcolorbox{bestpractice}[1][]{
  enhanced,breakable,
  colback=purple!5!white,colframe=purple!60!black,
  fonttitle=\bfseries,
  title=最佳实践,#1,
  boxed title style={colback=purple!60!black},
  attach boxed title to top left={yshift=-2mm,xshift=2mm},
  arc=2mm,boxrule=0.8mm,
}

% ==================== 章节标题 ====================
\usepackage{titlesec}
\usepackage{fancyhdr}
\usepackage{graphicx}
\usepackage{amssymb}
\usepackage{amsmath}
\usepackage{tabularx}
\usepackage{booktabs}
\usepackage{longtable}
\usepackage{array}
\usepackage{multirow}
\usepackage{enumitem}
\usepackage{tikz}
\usepackage{seqsplit}

\titleformat{\chapter}[display]
  {\normalfont\huge\bfseries\color{algheadbg}}
  {\chaptertitlename\ \thechapter}{20pt}{\Huge}
\titleformat{\section}
  {\normalfont\Large\bfseries\color{blue!70!black}}
  {\thesection}{1em}{}
\titleformat{\subsection}
  {\normalfont\large\bfseries\color{blue!60!black}}
  {\thesubsection}{1em}{}

\pagestyle{fancy}
\fancyhf{}
\fancyhead[LE,RO]{\thepage}
\fancyhead[RE]{\leftmark}
\fancyhead[LO]{\rightmark}

\setlist[itemize]{leftmargin=2em,itemsep=2pt}
\setlist[enumerate]{leftmargin=2em,itemsep=2pt}

% ==================== 超链接 ====================
\usepackage{hyperref}
\hypersetup{
  colorlinks=true,
  linkcolor=blue!70!black,
  urlcolor=blue!60!black,
  bookmarks=true,
  pdfauthor={Build Systems Reference},
  pdftitle={C/C++ Build Systems: CMake, XMake, VS, Xcode, Makefile},
}

% ==================== 自定义命令 ====================
\newcommand{\cpp}[1]{C++#1}
\newcommand{\Fn}[1]{\texttt{#1()}}
\newcommand{\Tp}[1]{\texttt{#1}}

% 语言标签
\newcommand{\cmakelabel}{%
  \noindent\begin{tcolorbox}[sharp corners,colback=cmakeheadbg,colframe=cmakeheadbg,
    left=6pt,right=6pt,top=2pt,bottom=2pt,boxrule=0pt]
  \textcolor{white}{\small\bfseries\sffamily CMake}\end{tcolorbox}%
}
\newcommand{\xmakelabel}{%
  \noindent\begin{tcolorbox}[sharp corners,colback=xmakeheadbg,colframe=xmakeheadbg,
    left=6pt,right=6pt,top=2pt,bottom=2pt,boxrule=0pt]
  \textcolor{white}{\small\bfseries\sffamily XMake}\end{tcolorbox}%
}
\newcommand{\vslabel}{%
  \noindent\begin{tcolorbox}[sharp corners,colback=vsheadbg,colframe=vsheadbg,
    left=6pt,right=6pt,top=2pt,bottom=2pt,boxrule=0pt]
  \textcolor{white}{\small\bfseries\sffamily Visual Studio / MSBuild}\end{tcolorbox}%
}
\newcommand{\xcodelabel}{%
  \noindent\begin{tcolorbox}[sharp corners,colback=xcodeheadbg,colframe=xcodeheadbg,
    left=6pt,right=6pt,top=2pt,bottom=2pt,boxrule=0pt]
  \textcolor{white}{\small\bfseries\sffamily Xcode}\end{tcolorbox}%
}
\newcommand{\makelabelcmd}{%
  \noindent\begin{tcolorbox}[sharp corners,colback=makeheadbg,colframe=makeheadbg,
    left=6pt,right=6pt,top=2pt,bottom=2pt,boxrule=0pt]
  \textcolor{white}{\small\bfseries\sffamily Makefile}\end{tcolorbox}%
}
\newcommand{\bashlabel}{%
  \noindent\begin{tcolorbox}[sharp corners,colback=cfgheadbg,colframe=cfgheadbg,
    left=6pt,right=6pt,top=2pt,bottom=2pt,boxrule=0pt]
  \textcolor{white}{\small\bfseries\sffamily Shell}\end{tcolorbox}%
}
\newcommand{\androidlabel}{%
  \noindent\begin{tcolorbox}[sharp corners,colback=androidheadbg!80!black,colframe=androidheadbg!80!black,
    left=6pt,right=6pt,top=2pt,bottom=2pt,boxrule=0pt]
  \textcolor{white}{\small\bfseries\sffamily Android}\end{tcolorbox}%
}
\newcommand{\oelabel}{%
  \noindent\begin{tcolorbox}[sharp corners,colback=oeheadbg,colframe=oeheadbg,
    left=6pt,right=6pt,top=2pt,bottom=2pt,boxrule=0pt]
  \textcolor{white}{\small\bfseries\sffamily OpenEmbedded / BitBake}\end{tcolorbox}%
}
\newcommand{\cfglabel}{%
  \noindent\begin{tcolorbox}[sharp corners,colback=cfgheadbg,colframe=cfgheadbg,
    left=6pt,right=6pt,top=2pt,bottom=2pt,boxrule=0pt]
  \textcolor{white}{\small\bfseries\sffamily Config}\end{tcolorbox}%
}

\makeindex

% ====================================================================
\begin{document}

% ==================== 封面 ====================
\frontmatter

\begin{titlepage}
\centering
\vspace*{3cm}
{\Huge\bfseries\color{algheadbg} C/C++ 构建系统参考手册\par}
\vspace{1cm}
{\LARGE\color{blue!70!black} CMake \textbullet{} XMake \textbullet{} Visual Studio \textbullet{} Xcode \textbullet{} Makefile\par}
\vspace{0.8cm}
{\Large\color{gray} 附录：Android 构建系统 \textbullet{} OpenEmbedded\par}
\vfill
{\large\color{gray} \today\par}
\end{titlepage}

\chapter*{前言}
\addcontentsline{toc}{chapter}{前言}

C 和 C++ 的构建系统生态丰富而复杂。从最传统的 Makefile，到工业标准的 CMake，再到新兴的 XMake，以及各大 IDE 各自的项目格式（Visual Studio 的 MSBuild、Xcode 的 pbxproj），开发者在不同场景下面临着不同的选择。嵌入式与移动端还有各自专属的构建系统——Android 的 Gradle + NDK 与 AOSP 构建、OpenEmbedded 的 BitBake 配方。

本文档系统梳理这些构建系统的核心概念、配置语法和最佳实践，帮助开发者在项目中选择合适的构建方案。每个构建系统都配有完整的配置示例和对比说明，附录部分介绍 Android 和 OpenEmbedded 这两个在嵌入式与移动开发中广泛使用的构建系统。

% --- TOC：从 build.ps1 生成的 manual.toc 读取目录 ---
\makeatletter
\newif\ifhavemanualtoc
\IfFileExists{\jobname-manual.toc}{\havemanualtoctrue}{}
\ifhavemanualtoc
  \clearpage
  \chapter*{\contentsname}%
  \@mkboth{\contentsname}{\contentsname}%
  \@input{\jobname-manual.toc}%
  \clearpage
\else
  \tableofcontents
\fi
\makeatother

% ==================== 正文 ====================
\mainmatter

% ####################################################################
\part{CMake}
% ####################################################################

% -------------------- Chapter 1 --------------------
\chapter{CMake 基础}

CMake（Cross-platform Make）是目前最广泛使用的跨平台构建系统。它本身不编译代码，而是根据 \Tp{CMakeLists.txt} 配置文件生成原生构建系统的文件（如 Unix 上的 Makefile、Windows 上的 Visual Studio 工程、macOS 上的 Xcode 工程），然后由原生工具完成实际编译。

\section{最小项目}

一个最简单的 CMake 项目只需要一个 \Tp{CMakeLists.txt} 文件：

\cmakelabel
\begin{lstlisting}[style=cmakestyle]
cmake_minimum_required(VERSION 3.20)
project(HelloWorld VERSION 1.0.0 LANGUAGES CXX)

add_executable(hello main.cpp)
\end{lstlisting}

\noindent
三行代码完成了三件事：指定最低 CMake 版本、声明项目名称和支持的语言、定义一个可执行文件目标。

\begin{guideline}[title=CMake 版本选择]
\begin{itemize}
  \item \textbf{3.15+}：支持 \Tp{CMAKE\_MSVC\_RUNTIME\_LIBRARY}，Windows 上可控制 CRT 链接方式。
  \item \textbf{3.20+}：支持 \Fn{cmake\_path}、\Tp{CMAKE\_CXX\_STANDARD} 设为 23。
  \item \textbf{3.24+}：\Tp{CMAKE\_COMPILE\_WARNING\_AS\_ERROR}，\Tp{FETCHCONTENT\_BASE\_DIR} 改进。
  \item \textbf{3.28+}：\Tp{CMAKE\_CXX\_STANDARD} 设为 26，\Tp{CMakePresets.json} v8。
\end{itemize}
\end{guideline}

\section{构建类型}

CMake 支持四种标准构建类型，通过 \Tp{CMAKE\_BUILD\_TYPE} 控制（仅对单配置生成器如 Makefile/Ninja 有效）：

\cmakelabel
\begin{lstlisting}[style=cmakestyle]
# 在命令行指定
# cmake -B build -DCMAKE_BUILD_TYPE=Release

# 或在 CMakeLists.txt 中设置默认值
if(NOT CMAKE_BUILD_TYPE AND NOT CMAKE_CONFIGURATION_TYPES)
  set(CMAKE_BUILD_TYPE Release CACHE STRING "Build type" FORCE)
  set_property(CACHE CMAKE_BUILD_TYPE PROPERTY
    STRINGS "Debug" "Release" "MinSizeRel" "RelWithDebInfo")
endif()
\end{lstlisting}

\begin{note}
Visual Studio 和 Xcode 等多配置生成器忽略 \Tp{CMAKE\_BUILD\_TYPE}，它们在构建时选择配置。可以用 \Tp{CMAKE\_CONFIGURATION\_TYPES} 限制可用配置。
\end{note}

\section{添加库}

CMake 支持静态库、共享库和接口库三种类型：

\cmakelabel
\begin{lstlisting}[style=cmakestyle]
# 静态库
add_library(mylib STATIC
  src/mylib.cpp
  src/utils.cpp
)

# 共享库（DLL / .so / .dylib）
add_library(myshared SHARED src/shared.cpp)

# 接口库（仅头文件，不生成二进制文件）
add_library(myheader_only INTERFACE)
target_include_directories(myheader_only
  INTERFACE ${CMAKE_CURRENT_SOURCE_DIR}/include)

# 推荐：使用 BUILD_SHARED_LIBS 变量控制库类型
option(BUILD_SHARED_LIBS "Build shared libraries" OFF)
add_library(mylib src/mylib.cpp)  (*@\# 类型由 BUILD\_SHARED\_LIBS 决定@*)
\end{lstlisting}

\section{Modern CMake：目标属性}

{\sloppy Modern CMake 的核心理念是 \textbf{一切皆目标}。避免使用目录级命令（如 \Fn{include\_directories}），改为目标级命令：\par}

\cmakelabel
\begin{lstlisting}[style=cmakestyle]
add_library(mylib src/mylib.cpp)

# 头文件路径 —— PRIVATE: 仅自身使用; INTERFACE: 仅传递给使用者
target_include_directories(mylib
  PRIVATE src
  INTERFACE include            (*@\# 使用者自动继承@*)
)

# 链接依赖
target_link_libraries(mylib PRIVATE fmt::fmt)

# C++ 标准
target_compile_features(mylib PUBLIC cxx_std_20)

# 编译定义
target_compile_definitions(mylib
  PRIVATE MYLIB_VERSION="1.0"
  PUBLIC $<BUILD_INTERFACE:MYLIB_BUILDING>
)
\end{lstlisting}

\begin{bestpractice}[title=PRIVATE / INTERFACE / PUBLIC 的语义]
\begin{itemize}
  \item \textbf{PRIVATE}：仅在构建此目标时使用，不传递给依赖者。
  \item \textbf{INTERFACE}：此目标自身不使用，但依赖者会使用。
  \item \textbf{PUBLIC}：自身和依赖者都会使用（= PRIVATE + INTERFACE）。
  \item 头文件库（INTERFACE library）只能用 INTERFACE 关键字。
\end{itemize}
\end{bestpractice}

\section{C++ 标准与编译选项}

\cmakelabel
\begin{lstlisting}[style=cmakestyle]
# 全局设置（不推荐用于大型项目）
set(CMAKE_CXX_STANDARD 20)
set(CMAKE_CXX_STANDARD_REQUIRED ON)
set(CMAKE_CXX_EXTENSIONS OFF)  (*@\# 禁用编译器特定扩展@*)

# 目标级设置（推荐）
target_compile_features(mylib PUBLIC cxx_std_20)

# 条件编译选项
if(MSVC)
  target_compile_options(mylib PRIVATE /W4 /permissive-)
else()
  target_compile_options(mylib PRIVATE -Wall -Wextra -Wpedantic)
endif()
\end{lstlisting}

\begin{warning}
\Tp{CMAKE\_CXX\_STANDARD} 是下限，不是精确版本。如果一个目标要求 \Tp{cxx\_std\_20}，另一个要求 \Tp{cxx\_std\_17}，CMake 会选择较高的那个。不要依赖它来降级标准。
\end{warning}

% -------------------- Chapter 2 --------------------
\chapter{CMake 进阶}

\section{find\_package 与依赖管理}

\Fn{find\_package} 是 CMake 查找第三方库的核心命令。它有两种模式：

\cmakelabel
\begin{lstlisting}[style=cmakestyle]
# Module 模式：查找 FindXXX.cmake 模块
find_package(Threads REQUIRED)

# Config 模式：查找 XXXConfig.cmake 或 xxx-config.cmake
find_package(fmt 10.0 REQUIRED CONFIG)
find_package(Boost 1.80 COMPONENTS filesystem system REQUIRED)

# 使用导入目标（imported target）
target_link_libraries(myapp PRIVATE
  Threads::Threads
  fmt::fmt
  Boost::filesystem
)
\end{lstlisting}

\begin{note}
\Fn{find\_package} 的搜索路径受 \Tp{CMAKE\_PREFIX\_PATH} 影响。安装库时指定 \Tp{-DCMAKE\_INSTALL\_PREFIX=/path}，使用者将同一路径加入 \Tp{CMAKE\_PREFIX\_PATH} 即可发现。
\end{note}

\section{FetchContent：下载时依赖}

从 CMake 3.11 起，\Tp{FetchContent} 模块可以在配置阶段自动下载并集成第三方库：

\cmakelabel
\begin{lstlisting}[style=cmakestyle]
include(FetchContent)

FetchContent_Declare(
  googletest
  GIT_REPOSITORY https://github.com/google/googletest.git
  GIT_TAG        v1.14.0
  GIT_SHALLOW    TRUE
)
# CMake 3.24+: 跳过测试和安装的构建
set(INSTALL_GTEST OFF CACHE BOOL "" FORCE)
set(gtest_force_shared_crt ON CACHE BOOL "" FORCE)

FetchContent_MakeAvailable(googletest)

# 现在可以使用 gtest 和 gtest_main 目标
add_executable(tests test_main.cpp)
target_link_libraries(tests PRIVATE gtest_main mylib)
\end{lstlisting}

\section{多目录项目结构}

大型项目通常将源码分散在多个子目录中，每个子目录有自己的 \Tp{CMakeLists.txt}：

\cmakelabel
\begin{lstlisting}[style=cmakestyle]
# 顶层 CMakeLists.txt
cmake_minimum_required(VERSION 3.20)
project(MyProject VERSION 2.0 LANGUAGES CXX)

set(CMAKE_CXX_STANDARD 20)
set(CMAKE_EXPORT_COMPILE_COMMANDS ON)  (*@\# 生成 compile\_commands.json@*)

add_subdirectory(external)   (*@\# 第三方依赖@*)
add_subdirectory(src/core)
add_subdirectory(src/app)
add_subdirectory(tests)
\end{lstlisting}

\begin{lstlisting}[style=cmakestyle]
# src/core/CMakeLists.txt
add_library(core STATIC core.cpp engine.cpp)
target_include_directories(core PUBLIC ${CMAKE_CURRENT_SOURCE_DIR})
target_link_libraries(core PRIVATE fmt::fmt)

# src/app/CMakeLists.txt
add_executable(myapp main.cpp)
target_link_libraries(myapp PRIVATE core)
\end{lstlisting}

\section{安装与导出}

为了让其他 CMake 项目可以通过 \Fn{find\_package} 找到你的库：

\cmakelabel
\begin{lstlisting}[style=cmakestyle]
include(GNUInstallDirs)

install(TARGETS mylib
  EXPORT mylibTargets
  ARCHIVE DESTINATION ${CMAKE_INSTALL_LIBDIR}
  LIBRARY DESTINATION ${CMAKE_INSTALL_LIBDIR}
  RUNTIME DESTINATION ${CMAKE_INSTALL_BINDIR}
  INCLUDES DESTINATION ${CMAKE_INSTALL_INCLUDEDIR}
)

install(DIRECTORY include/ DESTINATION ${CMAKE_INSTALL_INCLUDEDIR})

install(EXPORT mylibTargets
  FILE mylibTargets.cmake
  NAMESPACE mylib::
  DESTINATION ${CMAKE_INSTALL_LIBDIR}/cmake/mylib
)

# 生成 Config 文件
include(CMakePackageConfigHelpers)
configure_package_config_file(
  cmake/mylibConfig.cmake.in
  ${CMAKE_CURRENT_BINARY_DIR}/mylibConfig.cmake
  INSTALL_DESTINATION ${CMAKE_INSTALL_LIBDIR}/cmake/mylib
)
write_basic_package_version_file(
  ${CMAKE_CURRENT_BINARY_DIR}/mylibConfigVersion.cmake
  COMPATIBILITY SameMajorVersion
)
\end{lstlisting}

\section{CTest 测试框架}

\cmakelabel
\begin{lstlisting}[style=cmakestyle]
enable_testing()

add_executable(unit_tests test_core.cpp test_utils.cpp)
target_link_libraries(unit_tests PRIVATE core gtest_main)

include(GoogleTest)
gtest_discover_tests(unit_tests)

# 手动添加测试
add_test(NAME integration_test
  COMMAND ${CMAKE_COMMAND} -E env
    DATA_DIR=${CMAKE_SOURCE_DIR}/test_data
    $<TARGET_FILE:integration_test>
)
\end{lstlisting}

\section{Generator Expressions}

Generator Expressions 在生成阶段求值，可以实现条件化的编译选项：

\cmakelabel
\begin{lstlisting}[style=cmakestyle]
target_compile_definitions(mylib PRIVATE
  $<$<CONFIG:Debug>:DEBUG_MODE>
  $<$<CXX_COMPILER_ID:GNU>:GNU_SPECIFIC>
  $<$<PLATFORM_ID:Windows>:WIN32_LEAN_AND_MEAN>
)

# 仅在 Debug 配置下添加 ASan
target_compile_options(mylib PRIVATE
  $<$<CONFIG:Debug>:-fsanitize=address>
)
target_link_options(mylib PRIVATE
  $<$<CONFIG:Debug>:-fsanitize=address>
)
\end{lstlisting}

\section{CMakePresets.json}

\Tp{CMakePresets.json} 是 CMake 3.19 引入的标准化预设文件，统一管理配置、构建和测试参数：

\cfglabel
\begin{lstlisting}[style=bashstyle]
{
  "version": 6,
  "configurePresets": [
    {
      "name": "default",
      "generator": "Ninja",
      "binaryDir": "${sourceDir}/build",
      "cacheVariables": {
        "CMAKE_BUILD_TYPE": "Debug",
        "CMAKE_CXX_STANDARD": "20",
        "CMAKE_EXPORT_COMPILE_COMMANDS": "ON"
      }
    },
    {
      "name": "release",
      "inherits": "default",
      "cacheVariables": {
        "CMAKE_BUILD_TYPE": "Release"
      }
    }
  ],
  "buildPresets": [
    {
      "name": "default",
      "configurePreset": "default",
      "jobs": 8
    }
  ],
  "testPresets": [
    {
      "name": "default",
      "configurePreset": "default",
      "output": { "outputOnFailure": true }
    }
  ]
}
\end{lstlisting}

\begin{bestpractice}[title=CMake 项目最佳实践]
\begin{enumerate}
  \item 使用目标级命令（\Fn{target\_*}），避免目录级命令（\Fn{include\_directories} 等）。
  \item 设置 \Tp{CMAKE\_EXPORT\_COMPILE\_COMMANDS} 以生成 \Tp{compile\_commands.json}，供 clangd 等工具使用。
  \item 使用 \Tp{CMakePresets.json} 而非手写构建脚本来标准化配置。
  \item 第三方库优先使用 \Fn{find\_package}，不可用时使用 \Tp{FetchContent}。
  \item 始终使用命名空间化的导入目标（如 \Tp{fmt::fmt}）。
\end{enumerate}
\end{bestpractice}

% ####################################################################
\part{XMake}
% ####################################################################

\chapter{XMake 基础}

XMake 是一个基于 Lua 的轻量级跨平台构建工具。它使用 \Tp{xmake.lua} 作为配置文件，内置包管理器，语法比 CMake 更加简洁直观。

\section{最小项目}

\xmakelabel
\begin{lstlisting}[style=xmakestyle]
-- xmake.lua
add_rules("mode.debug", "mode.release")

target("hello")
    set_kind("binary")
    add_files("src/*.cpp")
\end{lstlisting}

\noindent
仅需四行即可完成项目配置。XMake 会自动检测系统上的编译器，无需手动指定工具链。

\section{构建类型}

\begin{lstlisting}[style=xmakestyle]
-- 内置的 debug/release 规则
add_rules("mode.debug", "mode.release")

-- 命令行切换
-- xmake f -m release
-- xmake
\end{lstlisting}

\noindent
\Tp{mode.debug} 规则自动添加调试标志和禁用优化；\Tp{mode.release} 启用优化并关闭调试信息。

\section{目标类型}

\begin{lstlisting}[style=xmakestyle]
-- 可执行文件
target("myapp")
    set_kind("binary")
    add_files("src/main.cpp")

-- 静态库
target("mylib")
    set_kind("static")
    add_files("src/lib/*.cpp")

-- 动态库
target("myshared")
    set_kind("shared")
    add_files("src/shared/*.cpp")

-- 头文件库（仅导出接口）
target("myheader")
    set_kind("headeronly")
    add_headerfiles("include/(**.h)")
\end{lstlisting}

\section{依赖与链接}

\begin{lstlisting}[style=xmakestyle]
target("myapp")
    set_kind("binary")
    add_files("src/*.cpp")
    add_deps("mylib")                     (*@\# 依赖项目内其他 target@*)
    add_links("pthread", "m")             (*@\# 链接系统库@*)
    add_includedirs("include")            (*@\# 头文件路径@*)
    add_linkdirs("/usr/local/lib")        (*@\# 库搜索路径@*)
    add_defines("USE_FEATURE")            (*@\# 宏定义@*)
\end{lstlisting}

\section{C++ 标准}

\begin{lstlisting}[style=xmakestyle]
target("myapp")
    set_kind("binary")
    add_files("src/*.cpp")
    set_languages("cxx20")               (*@\# C++20 标准@*)

-- 也可设置 C 标准
target("myclib")
    set_kind("static")
    add_files("src/*.c")
    set_languages("c17")                 (*@\# C17 标准@*)
\end{lstlisting}

\section{编译选项}

\begin{lstlisting}[style=xmakestyle]
target("myapp")
    set_kind("binary")
    add_files("src/*.cpp")
    set_languages("cxx20")

    -- 警告级别
    set_warnings("all", "extra", "error")

    -- 优化级别
    set_optimize("fastest")              -- none/fast/faster/fastest/smallest/aggressive

    -- 平台特定的编译标志
    if is_plat("windows") then
        add_cxflags("/W4", "/permissive-")
    else
        add_cxflags("-Wall", "-Wextra", "-Wpedantic")
    end
\end{lstlisting}

% -------------------- Chapter: XMake 进阶 --------------------
\chapter{XMake 进阶}

\section{包管理}

XMake 内置了远程包管理器（xrepo），可以自动下载和集成第三方库：

\begin{lstlisting}[style=xmakestyle]
-- 声明项目级依赖
add_requires("fmt", "zlib", "gtest")

target("myapp")
    set_kind("binary")
    add_files("src/*.cpp")
    add_packages("fmt", "zlib")

target("tests")
    set_kind("binary")
    add_files("tests/*.cpp")
    add_packages("gtest", "fmt")

-- 指定版本和配置
add_requires("boost >=1.80", {configs = {filesystem = true, system = true}})
add_requires("opencv 4.x", {configs = {cuda = true}})
\end{lstlisting}

\begin{lstlisting}[style=xmakestyle]
-- 添加自定义包源
add_repositories("myrepo https://github.com/user/xmake-repo.git")

-- 私有包：指定 URL 和哈希
add_requires("mylib", {
    urls = {"https://github.com/user/mylib/archive/v1.0.tar.gz"},
    configs = {shared = true}
})
\end{lstlisting}

\section{交叉编译}

XMake 内置了对交叉编译的支持，无需复杂的工具链文件：

\begin{lstlisting}[style=xmakestyle]
-- 命令行指定目标平台
-- xmake f -p android --ndk=/path/to/ndk
-- xmake f -p iphoneos -a arm64
-- xmake f -p mingw --sdk=/path/to/mingw

-- 在 xmake.lua 中配置工具链
set_toolchains("clang")                  -- 使用 clang

-- 自定义工具链
toolchain("my_arm_gcc")
    set_toolset("cc", "arm-linux-gnueabihf-gcc")
    set_toolset("cxx", "arm-linux-gnueabihf-g++")
    set_toolset("ld", "arm-linux-gnueabihf-g++")
    add_includedirs("/usr/arm-linux-gnueabihf/include")
    add_linkdirs("/usr/arm-linux-gnueabihf/lib")
toolchain_end()
\end{lstlisting}

\section{自定义规则}

\begin{lstlisting}[style=xmakestyle]
-- 自定义构建规则：Protocol Buffers
rule("protobuf")
    set_extensions(".proto")
    on_buildcmd_file(function (target, batchcmds, sourcefile, opt)
        local basename = path.basename(sourcefile)
        local outputdir = target:autogendir({root = true})
        local ccfile = path.join(outputdir, basename .. ".pb.cc")
        local hfile = path.join(outputdir, basename .. ".pb.h")

        batchcmds:show_progress(opt.progress,
            "compiling.proto %s", sourcefile)
        batchcmds:vrunv("protoc", {
            "--cpp_out=" .. outputdir, sourcefile
        })
        batchcmds:add_depfiles(sourcefile)
        batchcmds:set_depmtime(os.mtime(ccfile))
        batchcmds:set_depcache(target:dependfile(ccfile))
    end)
rule_end()

target("myapp")
    set_kind("binary")
    add_rules("protobuf")
    add_files("src/*.cpp")
    add_files("proto/*.proto")
\end{lstlisting}

\section{测试集成}

\begin{lstlisting}[style=xmakestyle]
add_requires("gtest")

target("tests")
    set_kind("binary")
    set_default(false)                   (*@\# 默认不构建@*)
    add_files("tests/*.cpp")
    add_packages("gtest")
    add_deps("mylib")

    -- 运行测试: xmake run tests
    on_run(function (target)
        os.exec(target:targetfile())
    end)
\end{lstlisting}

\begin{compare}[title=CMake vs XMake]
\begin{tabularx}{\textwidth}{lXX}
\toprule
\textbf{特性} & \textbf{CMake} & \textbf{XMake} \\
\midrule
配置语言 & CMake DSL & Lua \\
包管理 & 无内置（需 vcpkg/conan） & 内置 xrepo \\
学习曲线 & 中高 & 较低 \\
IDE 集成 & VS/CLion/VS Code 原生 & 需插件 \\
构建速度 & 中等 & 较快（内置缓存） \\
交叉编译 & 需工具链文件 & 内置支持 \\
社区规模 & 极大 & 增长中 \\
\bottomrule
\end{tabularx}
\end{compare}

% ####################################################################
\part{IDE 与系统构建工具}
% ####################################################################

\chapter{Visual Studio 项目系统}

Visual Studio 使用 MSBuild 作为底层构建引擎，项目文件由 \Tp{.vcxproj}（XML）和 \Tp{.sln}（解决方案）组成。虽然 CMake 可以直接生成 VS 工程，理解 MSBuild 的内部机制对于调试构建问题和定制构建流程仍然非常重要。

\section{解决方案与项目文件}

\Tp{.sln} 文件是一个纯文本格式，定义了解决方案中包含哪些项目以及构建配置的映射：

\vslabel
\begin{lstlisting}[style=bashstyle]
# .sln 文件结构（简化）
#
# Microsoft Visual Studio Solution File, Format Version 12.00
# Visual Studio Version 17
#
# Project("{GUID}") = "MyApp", "MyApp\MyApp.vcxproj", "{GUID}"
# EndProject
#
# GlobalSection(SolutionConfigurationPlatforms)
#   Debug|x64 = Debug|x64
#   Release|x64 = Release|x64
# EndGlobalSection
\end{lstlisting}

\section{vcxproj 结构}

\Tp{.vcxproj} 是一个 MSBuild XML 文件，定义了编译设置、源文件和依赖关系：

\vslabel
\begin{lstlisting}[style=xmlstyle]
<?xml version="1.0" encoding="utf-8"?>
<Project DefaultTargets="Build"
         ToolsVersion="17.0"
         xmlns="http://schemas.microsoft.com/developer/msbuild/2003">

  <!-- 全局属性 -->
  <PropertyGroup Label="Globals">
    <ProjectGuid>{12345678-1234-1234-1234-123456789ABC}</ProjectGuid>
    <RootNamespace>MyApp</RootNamespace>
    <WindowsTargetPlatformVersion>10.0</WindowsTargetPlatformVersion>
  </PropertyGroup>

  <!-- 导入默认属性 -->
  <Import Project="$(VCTargetsPath)\Microsoft.Cpp.Default.props" />

  <!-- Debug 配置 -->
  <PropertyGroup Label="Configuration"
                 Condition="'$(Configuration)|$(Platform)'=='Debug|x64'">
    <ConfigurationType>Application</ConfigurationType>
    <PlatformToolset>v143</PlatformToolset>
    <CharacterSet>Unicode</CharacterSet>
  </PropertyGroup>

  <!-- Release 配置 -->
  <PropertyGroup Label="Configuration"
                 Condition="'$(Configuration)|$(Platform)'=='Release|x64'">
    <ConfigurationType>Application</ConfigurationType>
    <PlatformToolset>v143</PlatformToolset>
    <CharacterSet>Unicode</CharacterSet>
    <WholeProgramOptimization>true</WholeProgramOptimization>
  </PropertyGroup>

  <Import Project="$(VCTargetsPath)\Microsoft.Cpp.props" />

  <!-- 编译选项 -->
  <ItemDefinitionGroup Condition="'$(Configuration)'=='Debug'">
    <ClCompile>
      <PreprocessorDefinitions>_DEBUG;%(PreprocessorDefinitions)
      </PreprocessorDefinitions>
      <RuntimeLibrary>MultiThreadedDebugDLL</RuntimeLibrary>
      <LanguageStandard>stdcpp20</LanguageStandard>
      <WarningLevel>Level4</WarningLevel>
      <AdditionalIncludeDirectories>include;%(AdditionalIncludeDirectories)
      </AdditionalIncludeDirectories>
    </ClCompile>
    <Link>
      <SubSystem>Console</SubSystem>
      <GenerateDebugInformation>true</GenerateDebugInformation>
    </Link>
  </ItemDefinitionGroup>

  <!-- 源文件 -->
  <ItemGroup>
    <ClCompile Include="src\main.cpp" />
    <ClCompile Include="src\app.cpp" />
    <ClInclude Include="include\app.h" />
  </ItemGroup>

  <Import Project="$(VCTargetsPath)\Microsoft.Cpp.targets" />
</Project>
\end{lstlisting}

\section{MSBuild 命令行}

\bashlabel
\begin{lstlisting}[style=bashstyle]
# 构建 Release|x64
msbuild MyApp.vcxproj /p:Configuration=Release /p:Platform=x64

# 清理并重建
msbuild MyApp.vcxproj /t:Rebuild /p:Configuration=Debug

# 并行编译
msbuild MyApp.sln /m:8 /p:Configuration=Release /p:Platform=x64

# 详细日志
msbuild MyApp.vcxproj /v:detailed /fl /flp:logfile=build.log
\end{lstlisting}

\section{NuGet 包管理}

\vslabel
\begin{lstlisting}[style=xmlstyle]
<!-- packages.config（传统方式） -->
<?xml version="1.0" encoding="utf-8"?>
<packages>
  <package id="fmt" version="10.2.1" targetFramework="native" />
  <package id="nlohmann.json" version="3.11.3" targetFramework="native" />
</packages>
\end{lstlisting}

\begin{lstlisting}[style=xmlstyle]
<!-- PackageReference（推荐方式，直接写在 vcxproj 中） -->
<ItemGroup>
  <PackageReference Include="fmt" Version="10.2.1" />
  <PackageReference Include="nlohmann.json" Version="3.11.3" />
</ItemGroup>
\end{lstlisting}

\section{属性表与目录文件}

属性表（\Tp{.props}）允许在多个项目之间共享构建设置：

\vslabel
\begin{lstlisting}[style=xmlstyle]
<!-- CommonSettings.props -->
<?xml version="1.0" encoding="utf-8"?>
<Project xmlns="http://schemas.microsoft.com/developer/msbuild/2003">
  <PropertyGroup>
    <LanguageStandard>stdcpp20</LanguageStandard>
    <ConformanceMode>true</ConformanceMode>
  </PropertyGroup>
  <ItemDefinitionGroup>
    <ClCompile>
      <AdditionalIncludeDirectories>$(SolutionDir)include;
        %(AdditionalIncludeDirectories)</AdditionalIncludeDirectories>
      <PreprocessorDefinitions>$(DefineConstants);
        %(PreprocessorDefinitions)</PreprocessorDefinitions>
    </ClCompile>
  </ItemDefinitionGroup>
</Project>
\end{lstlisting}

\begin{warning}
\Tp{Directory.Build.props} 和 \Tp{Directory.Build.targets} 文件会自动应用于同目录及子目录下的所有项目。这在大型解决方案中可以统一管理设置，但也可能导致难以追踪的隐式依赖。
\end{warning}

% -------------------- Chapter: Xcode --------------------
\chapter{Xcode 项目系统}

{\sloppy Xcode 使用 \Tp{.xcodeproj} 包（本质是一个目录）存储项目配置。其中核心文件是 \Tp{project.pbxproj}，一种类似旧式 plist 的文本格式。\par}

\section{Xcode 项目结构}

一个典型的 Xcode 项目目录结构如下：

\bashlabel
\begin{lstlisting}[style=bashstyle]
MyApp.xcodeproj/
  project.pbxproj              # 核心项目配置
  project.xcworkspace/        # 工作区配置
    contents.xcworkspacedata
  xcshareddata/
    xcschemes/                 # 共享 Scheme
      MyApp.xcscheme
  xcuserdata/
    user.xcuserdatad/          # 用户个人设置
\end{lstlisting}

\section{project.pbxproj 格式}

\Tp{pbxproj} 使用类似 NeXTSTEP plist 的语法。虽然通常由 Xcode IDE 管理，了解其结构对于手动修改和 CI 配置很有用：

\xcodelabel
\begin{lstlisting}[style=bashstyle]
// !$*UTF8*$!
{
  archiveVersion = 1;
  classes = {};
  objectVersion = 56;
  objects = {
    /* Begin PBXBuildFile section */
    A1B2C3 /* main.cpp in Sources */ = {
      isa = PBXBuildFile;
      fileRef = D4E5F6 /* main.cpp */;
    };
    /* End PBXBuildFile section */

    /* Begin PBXNativeTarget section */
    ABC123 /* MyApp */ = {
      isa = PBXNativeTarget;
      buildConfigurationList = DEF456;
      buildPhases = (
        GHI789 /* Sources */,
        JKL012 /* Frameworks */,
        MNO345 /* Resources */,
      );
      buildRules = ();
      dependencies = ();
      name = MyApp;
      productName = MyApp;
      productReference = PQR678 /* MyApp */;
      productType = "com.apple.product-type.tool";
    };
    /* End PBXNativeTarget section */

    /* Begin XCBuildConfiguration section */
    CFG001 /* Debug */ = {
      isa = XCBuildConfiguration;
      buildSettings = {
        CLANG_CXX_LANGUAGE_STANDARD = "c++20";
        CLANG_ENABLE_MODULES = YES;
        GCC_PREPROCESSOR_DEFINITIONS = ("DEBUG=1",);
        PRODUCT_NAME = "$(TARGET_NAME)";
        SDKROOT = macosx;
      };
      name = Debug;
    };
    CFG002 /* Release */ = {
      isa = XCBuildConfiguration;
      buildSettings = {
        CLANG_CXX_LANGUAGE_STANDARD = "c++20";
        GCC_OPTIMIZATION_LEVEL = s;
        PRODUCT_NAME = "$(TARGET_NAME)";
        SDKROOT = macosx;
      };
      name = Release;
    };
    /* End XCBuildConfiguration section */
  };
  rootObject = ROOT01 /* Project object */;
}
\end{lstlisting}

\section{xcodebuild 命令行}

\bashlabel
\begin{lstlisting}[style=bashstyle]
# 列出项目中的所有 scheme
xcodebuild -list -project MyApp.xcodeproj

# 构建指定 scheme
xcodebuild build \
  -project MyApp.xcodeproj \
  -scheme MyApp \
  -configuration Release \
  -derivedDataPath build/

# 构建并运行测试
xcodebuild test \
  -project MyApp.xcodeproj \
  -scheme MyAppTests \
  -destination 'platform=macOS'

# 清理
xcodebuild clean -project MyApp.xcodeproj -scheme MyApp

# 导出 archive
xcodebuild archive \
  -project MyApp.xcodeproj \
  -scheme MyApp \
  -archivePath build/MyApp.xcarchive
\end{lstlisting}

\section{Swift Package Manager 集成}

Xcode 11+ 原生支持 Swift Package Manager（SPM），可以直接在 \Tp{.xcodeproj} 中声明 C/C++ 依赖：

\xcodelabel
\begin{lstlisting}[style=bashstyle]
// Package.swift（C++ 库的 SPM 配置）
// swift-tools-version:5.9
import PackageDescription

let package = Package(
    name: "MyLib",
    products: [
        .library(name: "MyLib", targets: ["MyLib"]),
    ],
    targets: [
        .target(
            name: "MyLib",
            path: "Sources",
            publicHeadersPath: "include",
            cxxSettings: [
                .headerSearchPath("internal"),
                .define("MYLIB_BUILDING"),
                .unsafeFlags(["-std=c++20"]),
            ]
        ),
        .testTarget(
            name: "MyLibTests",
            dependencies: ["MyLib"]
        ),
    ],
    cxxLanguageStandard: .cxx20
)
\end{lstlisting}

\section{CMake 生成 Xcode 工程}

CMake 可以直接生成 Xcode 项目，这在跨平台项目中很常见：

\bashlabel
\begin{lstlisting}[style=bashstyle]
# 生成 Xcode 工程
cmake -B build -G Xcode -DCMAKE_CXX_STANDARD=20

# 打开
open build/MyProject.xcodeproj

# 或通过命令行构建
cmake --build build --config Release
\end{lstlisting}

\begin{compare}[title=VS 项目 vs Xcode 项目]
\begin{tabularx}{\textwidth}{lXX}
\toprule
\textbf{特性} & \textbf{Visual Studio} & \textbf{Xcode} \\
\midrule
项目文件格式 & XML (\Tp{.vcxproj}) & NeXTSTEP plist (\Tp{.pbxproj}) \\
解决方案 & \Tp{.sln} & \Tp{.xcworkspace} \\
构建引擎 & MSBuild & xcodebuild \\
配置管理 & Property Sheet (\Tp{.props}) & xcconfig 文件 \\
包管理 & NuGet / vcpkg & SPM / CocoaPods \\
命令行工具 & \Tp{msbuild} & \Tp{xcodebuild} \\
跨平台 & Windows / Linux (WSL) & macOS / iOS \\
\bottomrule
\end{tabularx}
\end{compare}

% -------------------- Chapter: Makefile --------------------
\chapter{Makefile 与 GNU Make}

Make 是最古老的构建工具之一，至今仍在 Unix 世界中广泛使用。GNU Make 是其最流行的实现，几乎所有 Linux 发行版都预装了它。虽然现代项目倾向于使用 CMake 或 XMake，但许多遗留项目和底层系统组件仍在使用 Makefile。

\section{基本规则}

Makefile 由一系列 \textbf{规则} 组成，每条规则定义了目标（target）、依赖（prerequisites）和命令（recipe）：

\makelabelcmd
\begin{lstlisting}[style=bashstyle]
# 目标: 依赖列表
# \t命令（必须是 Tab 缩进）

hello: main.o utils.o
\tg++ -o hello main.o utils.o

main.o: main.cpp utils.h
\tg++ -c main.cpp

utils.o: utils.cpp utils.h
\tg++ -c utils.cpp

clean:
\trm -f *.o hello
\end{lstlisting}

\begin{warning}
Makefile 中的命令必须使用 \textbf{Tab 字符}缩进，不能使用空格。这是 Make 最常被抱怨的设计缺陷之一。如果看到 ``missing separator'' 错误，几乎总是因为用了空格代替 Tab。
\end{warning}

\section{变量}

\makelabelcmd
\begin{lstlisting}[style=bashstyle]
# 变量定义
CXX      = g++
CXXFLAGS = -std=c++20 -Wall -Wextra
LDFLAGS  = -L/usr/local/lib
LDLIBS   = -lfmt -lpthread
SRCDIR   = src
OBJDIR   = build

# 递归展开（使用时求值）
SOURCES = $(wildcard $(SRCDIR)/*.cpp)

# 简单展开（定义时求值）
OBJECTS := $(patsubst $(SRCDIR)/%.cpp,$(OBJDIR)/%.o,$(SOURCES))

# 追加
CXXFLAGS += -O2

# 条件赋值（仅在未定义时赋值）
CXX ?= clang++
\end{lstlisting}

\section{隐式规则与模式规则}

\makelabelcmd
\begin{lstlisting}[style=bashstyle]
# 隐式规则：Make 内置了 .cpp -> .o 的规则
# 等价于：
%.o: %.cpp
\t$(CXX) $(CXXFLAGS) -c $< -o $@

# 模式规则：编译 src/ 下的 .cpp 到 build/ 下的 .o
$(OBJDIR)/%.o: $(SRCDIR)/%.cpp | $(OBJDIR)
\t$(CXX) $(CXXFLAGS) -c $< -o $@

# 目录创建（order-only prerequisite）
$(OBJDIR):
\tmkdir -p $(OBJDIR)

# 自动变量说明：
# $@ = 目标文件名
# $< = 第一个依赖
# $^ = 所有依赖（去重）
# $? = 比目标更新的依赖
\end{lstlisting}

\section{完整的 Makefile 示例}

\makelabelcmd
\begin{lstlisting}[style=bashstyle]
# ==================== 配置 ====================
CXX       := g++
CXXFLAGS  := -std=c++20 -Wall -Wextra -Wpedantic
LDFLAGS   :=
LDLIBS    := -lfmt -lpthread

SRCDIR    := src
BUILDDIR  := build
TARGET    := myapp

# ==================== 自动推导 ====================
SOURCES   := $(shell find $(SRCDIR) -name '*.cpp')
OBJECTS   := $(patsubst $(SRCDIR)/%.cpp,$(BUILDDIR)/%.o,$(SOURCES))
DEPS      := $(OBJECTS:.o=.d)

# ==================== 规则 ====================
.PHONY: all clean debug release

all: $(TARGET)

$(TARGET): $(OBJECTS)
\t$(CXX) $(LDFLAGS) -o $@ $^ $(LDLIBS)

$(BUILDDIR)/%.o: $(SRCDIR)/%.cpp | $(BUILDDIR)
\t@mkdir -p $(dir $@)
\t$(CXX) $(CXXFLAGS) -MMD -MP -c $< -o $@

$(BUILDDIR):
\t@mkdir -p $@

# 自动依赖追踪
-include $(DEPS)

clean:
\trm -rf $(BUILDDIR) $(TARGET)

debug: CXXFLAGS += -g -O0 -DDEBUG
debug: all

release: CXXFLAGS += -O2 -DNDEBUG
release: all
\end{lstlisting}

\section{高级特性}

\makelabelcmd
\begin{lstlisting}[style=bashstyle]
# 条件判断
ifeq ($(OS),Windows_NT)
  TARGET := $(TARGET).exe
  LDLIBS += -lws2_32
else
  UNAME_S := $(shell uname -s)
  ifeq ($(UNAME_S),Darwin)
    CXXFLAGS += -stdlib=libc++
  endif
endif

# 函数使用
COMMA := ,
EMPTY :=
SPACE := $(EMPTY) $(EMPTY)

# 字符串替换
UPPER_SOURCES := $(subst .cpp,.CPP,$(SOURCES))

# 循环
LIBS := fmt zlib boost_system
$(foreach lib,$(LIBS),$(eval $(lib)_CFLAGS := $(shell pkg-config --cflags $(lib))))

# include 其他 Makefile
-include config.mk           (*@\# - 前缀：文件不存在时不报错@*)
\end{lstlisting}

\section{Make 与 CMake 的关系}

在实际项目中，Make 和 CMake 并不矛盾。CMake 可以生成 Makefile 作为其底层构建系统：

\bashlabel
\begin{lstlisting}[style=bashstyle]
# CMake 生成 Makefile
cmake -B build -G "Unix Makefiles"
cmake --build build -j8

# 或直接使用生成的 Makefile
cd build && make -j8

# 也可以使用 Ninja（更快）
cmake -B build -G Ninja
cmake --build build
\end{lstlisting}

\begin{compare}[title=Make vs CMake vs XMake]
\begin{tabularx}{\textwidth}{lXXX}
\toprule
\textbf{特性} & \textbf{Make} & \textbf{CMake} & \textbf{XMake} \\
\midrule
跨平台 & Unix 为主 & 全平台 & 全平台 \\
配置语法 & Makefile DSL & CMake DSL & Lua \\
依赖管理 & 无 & 无内置 & 内置 xrepo \\
增量编译 & 手动管理 & 自动 & 自动 \\
并行构建 & \Tp{make -j} & \Tp{cmake --build -j} & \Tp{xmake -j} \\
学习曲线 & 中等 & 中高 & 较低 \\
\bottomrule
\end{tabularx}
\end{compare}

\begin{bestpractice}[title=Makefile 最佳实践]
\begin{enumerate}
  \item 使用 \Tp{-MMD -MP} 生成 \Tp{.d} 依赖文件，实现头文件变更自动重编译。
  \item 使用 \Tp{.PHONY} 声明伪目标（\Tp{clean}、\Tp{all} 等），避免与同名文件冲突。
  \item 使用 \Tp{@} 前缀隐藏命令输出，使用 \Tp{-} 前缀忽略错误。
  \item 使用 order-only prerequisite（\Tp{| dir}）确保目录先于文件创建。
  \item 目标特定变量（如 \Tp{debug: CXXFLAGS += -g}）可以为不同配置定制编译选项。
\end{enumerate}
\end{bestpractice}

% ####################################################################
\part{附录}
% ####################################################################
\appendix

% -------------------- Appendix A: Android --------------------
\chapter{Android 构建系统}

Android 开发涉及多层次的构建系统：应用层使用 Gradle，NDK 层使用 CMake 或 ndk-build，系统层（AOSP）使用 Soong 和 Make。

\section{Gradle + CMake（NDK）}

这是目前 Android 原生开发的标准方案。Gradle 作为顶层构建工具，通过 externalNativeBuild 调用 CMake 编译 C/C++ 代码：

\androidlabel
\begin{lstlisting}[style=bashstyle]
// build.gradle.kts（Kotlin DSL）
plugins {
    id("com.android.application") version "8.2.0"
}

android {
    namespace = "com.example.myapp"
    compileSdk = 34
    ndkVersion = "26.1.10909125"

    defaultConfig {
        applicationId = "com.example.myapp"
        minSdk = 24
        targetSdk = 34

        externalNativeBuild {
            cmake {
                cppFlags("-std=c++20", "-fno-rtti")
                arguments("-DANDROID_STL=c++_shared",
                          "-DCMAKE_BUILD_TYPE=Release")
                abiFilters("arm64-v8a", "x86_64")
            }
        }
    }

    externalNativeBuild {
        cmake {
            path("src/main/cpp/CMakeLists.txt")
            version = "3.22.1"
        }
    }
}
\end{lstlisting}

\cmakelabel
\begin{lstlisting}[style=cmakestyle]
# src/main/cpp/CMakeLists.txt
cmake_minimum_required(VERSION 3.22.1)
project(nativelib)

add_library(nativelib SHARED
  native-lib.cpp
  jni-bridge.cpp
)

# NDK 提供的变量
target_include_directories(nativelib PRIVATE
  ${ANDROID_NDK}/sources/android/native_app_glue
)

# 链接 Android 日志库
find_library(log-lib log)
target_link_libraries(nativelib
  ${log-lib}
  android           (*@\# NDK 提供的 android 库@*)
)

# C++20
target_compile_features(nativelib PRIVATE cxx_std_20)
\end{lstlisting}

\section{NDK 特有变量}

CMake 在 NDK 环境中会收到以下预定义变量：

\begin{guideline}[title=NDK CMake 变量]
\begin{tabularx}{\textwidth}{lX}
\toprule
\textbf{变量} & \textbf{说明} \\
\midrule
\Tp{ANDROID\_ABI} & 目标 ABI（\Tp{arm64-v8a}、\Tp{x86\_64} 等） \\
\Tp{ANDROID\_PLATFORM} & 最低 API 级别（如 \Tp{android-24}） \\
\Tp{ANDROID\_NDK} & NDK 根目录路径 \\
\Tp{ANDROID\_STL} & C++ 标准库选择（\Tp{c++\_shared} 或 \Tp{c++\_static}） \\
\Tp{CMAKE\_TOOLCHAIN\_FILE} & NDK 提供的工具链文件路径 \\
\Tp{CMAKE\_SYSTEM\_NAME} & 值为 \Tp{Android} \\
\bottomrule
\end{tabularx}
\end{guideline}

\section{ndk-build}

{\sloppy\Tp{ndk-build} 是 NDK 的旧式构建系统，使用 \Tp{Android.mk}（Makefile 方言）和 \Tp{Application.mk}：\par}

\androidlabel
\begin{lstlisting}[style=bashstyle]
# Android.mk
LOCAL_PATH := $(call my-dir)

include $(CLEAR_VARS)
LOCAL_MODULE    := nativelib
LOCAL_SRC_FILES := native-lib.cpp jni-bridge.cpp
LOCAL_C_INCLUDES := $(LOCAL_PATH)/include
LOCAL_LDLIBS    := -llog -landroid
LOCAL_CPP_FEATURES := rtti exceptions
LOCAL_CFLAGS    := -std=c++20

include $(BUILD_SHARED_LIBRARY)
\end{lstlisting}

\begin{lstlisting}[style=bashstyle]
# Application.mk
APP_STL     := c++_shared
APP_ABI     := arm64-v8a x86_64
APP_PLATFORM := android-24
APP_CPPFLAGS := -std=c++20
\end{lstlisting}

\begin{note}
Google 已宣布 ndk-build 进入维护模式，新项目应使用 CMake。但大量旧项目（包括 Android 系统源码中的部分模块）仍在使用 ndk-build。
\end{note}

\section{AOSP 构建系统：Soong + Make}

Android 开源项目（AOSP）使用 Soong（基于 Blueprint）和传统 Make 的混合构建系统：

\androidlabel
\begin{lstlisting}[style=bashstyle]
# Android.bp（Soong 配置，类似 JSON）
cc_library_shared {
    name: "libnativelib",
    srcs: ["native-lib.cpp", "jni-bridge.cpp"],
    shared_libs: [
        "liblog",
        "libandroid",
    ],
    cflags: [
        "-Wall",
        "-Werror",
    ],
    cppflags: ["-std=c++20"],
    stl: "libc++",
    sdk_version: "current",
}
\end{lstlisting}

\begin{lstlisting}[style=bashstyle]
# Android.mk（传统 Make 方式，仍然存在于部分模块中）
LOCAL_PATH := $(call my-dir)
include $(CLEAR_VARS)
LOCAL_MODULE := libnativelib
LOCAL_SRC_FILES := native-lib.cpp
LOCAL_SHARED_LIBRARIES := liblog libandroid
include $(BUILD_SHARED_LIBRARY)
\end{lstlisting}

\begin{bestpractice}[title=Android 构建系统选择]
\begin{itemize}
  \item \textbf{应用开发}：Gradle + CMake（NDK），这是官方推荐的方案。
  \item \textbf{纯 C/C++ 库}：独立 CMake 项目，通过 Gradle 集成。
  \item \textbf{AOSP 系统开发}：优先使用 \Tp{Android.bp}（Soong），避免新增 \Tp{Android.mk}。
  \item \textbf{跨平台库}：使用 CMake 作为统一构建系统，为 Android 添加 NDK 工具链文件。
\end{itemize}
\end{bestpractice}

% -------------------- Appendix B: OpenEmbedded --------------------
\chapter{OpenEmbedded 构建系统}

OpenEmbedded（OE）是一个用于构建嵌入式 Linux 发行版的元构建系统。它的核心工具是 \textbf{BitBake}——一个类似 Make 的任务执行引擎，通过 \textbf{配方}（recipe）文件定义如何获取、配置、编译和安装软件包。Yocto Project 是 OE 最著名的上层发行版。

\section{核心概念}

\begin{guideline}[title=OE 核心术语]
\begin{tabularx}{\textwidth}{lX}
\toprule
\textbf{术语} & \textbf{说明} \\
\midrule
\textbf{Recipe} (\Tp{.bb}) & 定义如何构建一个软件包 \\
\textbf{Layer} & 一组相关的 recipe 集合 \\
\textbf{BitBake} & 任务调度和执行引擎 \\
\textbf{Machine} & 目标硬件平台 \\
\textbf{Distro} & 目标 Linux 发行版配置 \\
\textbf{Image} & 最终的文件系统镜像 \\
\textbf{Sysroot} & 交叉编译的 sysroot 环境 \\
\textbf{WORKDIR} & BitBake 的临时工作目录 \\
\bottomrule
\end{tabularx}
\end{guideline}

\section{Recipe 文件结构}

一个典型的 BitBake recipe（\Tp{.bb} 文件）：

\oelabel
\begin{lstlisting}[style=bbstyle]
# myapp_1.0.bb
SUMMARY = "My C++ Application"
DESCRIPTION = "A sample C++20 application for embedded Linux"
HOMEPAGE = "https://github.com/user/myapp"
LICENSE = "MIT"
LIC_FILES_CHKSUM = "file://LICENSE;md5=abc123..."

# 源码获取
SRC_URI = "git://github.com/user/myapp.git;protocol=https;branch=main \
           file://0001-fix-build.patch"
SRCREV = "a1b2c3d4e5f6..."
S = "${WORKDIR}/git"

# 依赖
DEPENDS = "fmt zlib boost"
RDEPENDS:${PN} = "fmt zlib"

# 使用 CMake 构建
inherit cmake

# 额外的 CMake 参数
EXTRA_OECMAKE = "-DENABLE_FEATURE=ON -DBUILD_TESTS=OFF"

# 安装额外文件
do_install:append() {
    install -d ${D}${sysconfdir}
    install -m 0644 ${S}/config/myapp.conf ${D}${sysconfdir}/
}

# 包拆分
PACKAGES =+ "${PN}-config ${PN}-doc"
FILES:${PN}-config = "${sysconfdir}/myapp.conf"
FILES:${PN}-doc = "${docdir}"
\end{lstlisting}

\section{BitBake 任务流}

BitBake 的标准任务执行流程：

\bashlabel
\begin{lstlisting}[style=bashstyle]
# 构建流程（按顺序执行）
do_fetch      # 下载源码
do_unpack     # 解压到 WORKDIR
do_patch      # 应用补丁
do_configure  # 运行 configure（CMake: cmake）
do_compile    # 编译（CMake: make/ninja）
do_install    # 安装到 D（目标 sysroot）
do_package    # 生成二进制包（.deb/.rpm/.ipk）
do_populate_sysroot  # 填充 sysroot 供其他 recipe 使用
\end{lstlisting}

\section{继承与类}

Recipe 通过 \Tp{inherit} 复用构建逻辑。OE 提供了大量内置类：

\oelabel
\begin{lstlisting}[style=bbstyle]
# CMake 项目
inherit cmake
EXTRA_OECMAKE = "-DCMAKE_BUILD_TYPE=Release"

# Autotools 项目
inherit autotools
EXTRA_OECONF = "--enable-feature --disable-tests"

# Meson 项目
inherit meson
EXTRA_OEMESON = "-Dfeature=enabled"

# Python 项目
inherit setuptools3

# 纯脚本/配置文件（不需要编译）
inherit allarch

# 自定义任务
do_mytask() {
    echo "Running custom task..."
}
addtask mytask after do_compile before do_install
\end{lstlisting}

\section{Layer 管理}

\oelabel
\begin{lstlisting}[style=bashstyle]
# 添加自定义 layer
bitbake-layers create-layer meta-mycustom
bitbake-layers add-layer meta-mycustom

# layer 目录结构
meta-mycustom/
  conf/
    layer.conf              # Layer 配置
  recipes-myapp/
    myapp/
      myapp_1.0.bb          # Recipe 文件
      myapp/
        0001-fix-build.patch  # 补丁文件
        myapp.service        # systemd 服务文件
\end{lstlisting}

\cfglabel
\begin{lstlisting}[style=bashstyle]
# conf/layer.conf
BBPATH .= ":${LAYERDIR}"
BBFILE_COLLECTIONS += "meta-mycustom"
BBFILE_PATTERN_meta-mycustom = "^${LAYERDIR}/"
BBFILE_PRIORITY_meta-mycustom = "6"
LAYERSERIES_COMPAT_meta-mycustom = "scarthgap"
\end{lstlisting}

\section{镜像构建}

\oelabel
\begin{lstlisting}[style=bbstyle]
# my-image.bb - 自定义镜像
SUMMARY = "Custom embedded Linux image"

inherit core-image

IMAGE_INSTALL = " \
    packagegroup-core-boot \
    myapp \
    myapp-config \
    openssh \
    systemd \
    kernel-modules \
"

IMAGE_FEATURES += "ssh-server-openssh debug-tweaks"

# 镜像大小限制（MB）
IMAGE_ROOTFS_EXTRA_SPACE = "512000"
\end{lstlisting}

\bashlabel
\begin{lstlisting}[style=bashstyle]
# 构建镜像
bitbake my-image

# 构建 SDK（交叉编译工具链）
bitbake my-image -c populate_sdk

# DevShell（进入交叉编译环境）
bitbake myapp -c devshell

# 查看依赖图
bitbake -g my-image && cat task-depends.dot
\end{lstlisting}

\section{Append 与 bbappend}

\Tp{.bbappend} 文件允许在不修改原始 recipe 的情况下追加或覆盖设置：

\oelabel
\begin{lstlisting}[style=bbstyle]
# myapp_1.0.bbappend（在你的 layer 中）
# 文件名必须与原始 recipe 匹配（支持通配符）

# 添加额外的源码补丁
SRC_URI += "file://0002-custom-patch.patch"

# 修改编译选项
EXTRA_OECMAKE += "-DENABLE_EXTRA=ON"

# 追加安装步骤
do_install:append() {
    install -d ${D}${systemd_system_unitdir}
    install -m 0644 ${WORKDIR}/myapp.service \
        ${D}${systemd_system_unitdir}/
}

# 添加 systemd 服务
inherit systemd
SYSTEMD_SERVICE:${PN} = "myapp.service"
SYSTEMD_AUTO_ENABLE = "enable"
\end{lstlisting}

\begin{compare}[title=Android 构建 vs OpenEmbedded]
\begin{tabularx}{\textwidth}{lXX}
\toprule
\textbf{特性} & \textbf{Android (AOSP)} & \textbf{OpenEmbedded} \\
\midrule
构建引擎 & Soong + Make & BitBake \\
配置文件 & \Tp{Android.bp} / \Tp{Android.mk} & \Tp{.bb} / \Tp{.bbappend} \\
目标平台 & Android 设备 & 嵌入式 Linux \\
包管理 & 无（整体镜像） & \Tp{ipk} / \Tp{deb} / \Tp{rpm} \\
交叉编译 & NDK 工具链 & OE sysroot \\
层/模块化 & 模块 (\Tp{Android.bp}) & Layer (\Tp{meta-*}) \\
C++ 标准库 & libc++ (NDK) & 可配置 (glibc/musl) \\
\bottomrule
\end{tabularx}
\end{compare}

\begin{bestpractice}[title=OpenEmbedded 最佳实践]
\begin{enumerate}
  \item 使用 \Tp{.bbappend} 扩展已有 recipe，不要复制整个 recipe。
  \item 将自定义 recipe 组织到独立的 layer 中，保持 layer 职责单一。
  \item 使用 \Tp{SRCREV} 锁定 Git 提交哈希，避免 \Tp{AUTOREV}（自动获取最新版本）用于生产构建。
  \item 善用 \Tp{PACKAGECONFIG} 实现可选功能的开关管理。
  \item 使用 \Tp{sstate-cache} 加速重复构建，确保 CI 环境共享 sstate 目录。
\end{enumerate}
\end{bestpractice}

% ==================== 后记 ====================
\backmatter

\chapter*{构建系统总览}
\addcontentsline{toc}{chapter}{构建系统总览}

下表从多个维度对比了本文涵盖的所有构建系统：

{\small
\begin{longtable}{>{\bfseries}l p{2.5cm} p{2.5cm} p{2.5cm} p{2.5cm}}
\toprule
\textbf{构建系统} & \textbf{适用场景} & \textbf{配置语言} & \textbf{跨平台} & \textbf{包管理} \\
\midrule
CMake & 通用 C/C++ & CMake DSL & 全平台 & 无内置 \\
XMake & 通用 C/C++ & Lua & 全平台 & xrepo \\
MSBuild & Windows / VS & XML & Windows 为主 & NuGet / vcpkg \\
Xcode & macOS / iOS & pbxproj & Apple 平台 & SPM \\
Makefile & Unix 系统 & Make DSL & Unix 为主 & 无 \\
\midrule
Gradle+CMake & Android NDK & Kotlin/Groovy & Android & Maven \\
AOSP Soong & Android 系统 & Blueprint & Android & 模块系统 \\
BitBake & 嵌入式 Linux & BitBake DSL & 嵌入式 & ipk/deb/rpm \\
\bottomrule
\end{longtable}
}

选择构建系统时，需要考虑项目规模、目标平台、团队熟悉度和生态系统兼容性。对于新项目，CMake 仍然是最安全的选择——几乎所有 IDE 和 CI 系统都原生支持它。如果追求更简洁的语法和内置包管理，XMake 值得一试。对于特定平台（如 Android 或嵌入式 Linux），则需要配合使用平台专属的构建系统。

\end{document}
